\documentclass[journal]{IEEEtran}
\usepackage{amsmath,amssymb,amsfonts,amsthm}
\usepackage{booktabs}
\usepackage{hyperref}

\newtheorem{definition}{Definition}
\newtheorem{axiom}{Axiom}
\newtheorem{theorem}{Theorem}

\title{Causal Transparency Algebra ($\mathbb{R}_\emptyset$) and $\emptyset$-Net: A Non-Distributive Bi-Identity Framework for Drift-Free Continual Learning and Zero-Power Neuromorphic Bypass}
\author{Hakim Guettal}

\begin{document}
\maketitle

\begin{abstract}
We introduce Causal Transparency Algebra ($\mathbb{R}_\emptyset = \mathbb{R} \sqcup \{\emptyset\}$) and the $\emptyset$-Net neural architecture. Unlike IEEE 754 \texttt{NaN}, Codd's \texttt{NULL}, or Wheel Algebras, $\emptyset$ acts as a simultaneous bilateral additive and multiplicative identity ($x \oplus \emptyset = x$, $x \otimes \emptyset = x$) with quotient collapse ($x \oslash \emptyset = \emptyset$). We prove the Causal Arrow Theorem and derive the internal polynomial detector $\chi_\emptyset(x) = ((1 \otimes 0) \oplus (1 \otimes x)) \ominus (1 \otimes (0 \oplus x))$, which extracts hardware power-gating bits directly from controlled non-distributivity without conditional branches.
\end{abstract}

\section{Axioms of $\mathbb{R}_\emptyset$}
\begin{axiom}[Universal Bilateral Identity]
$\forall x \in \mathbb{R}_\emptyset, \; x \oplus \emptyset = \emptyset \oplus x = x$ and $x \otimes \emptyset = \emptyset \otimes x = x$.
\end{axiom}
\begin{axiom}[Quotient Collapse]
$\forall x \in \mathbb{R}_\emptyset, \; x \oslash \emptyset = \emptyset, \; \emptyset \oslash x = \emptyset, \; \emptyset \oslash 0 = \emptyset$.
\end{axiom}
\begin{axiom}[A-Polarity and Isolation]
$\ominus \emptyset = \emptyset$, and $\{\emptyset\}$ is clopen in $(\mathbb{R}_\emptyset, \tau_\emptyset)$.
\end{axiom}

\section{The Algebraic Transparency Detector $\chi_\emptyset(x)$}
\begin{theorem}
The polynomial functional $\chi_\emptyset(x) \triangleq ((1 \otimes 0) \oplus (1 \otimes x)) \ominus (1 \otimes (0 \oplus x))$ evaluates to $0$ for all $x \in \mathbb{R}$ (including $x=0$) and to $1$ for $x = \emptyset$.
\end{theorem}

\end{document}
